\section{Implementazione con esempio pratico}
Il processo di hydration è relativamente semplice a livello concettuale, ma può essere complesso da implementare in pratica. In questa sezione, vediamo nel dettaglio come funziona in un'applicazione web moderna, prendendo come esempio la schermata home di YouTube.

Come visto in precedenza, il processo di hydration è composto da due fasi: la prima fase è la renderizzazione del contenuto lato server, mentre la seconda fase è la vera e propria hydration del contenuto lato client.

\subsection{Prima fase: renderizzazione lato server}

\vspace{5mm}
\begin{figure}[H]
    \centering
    \includegraphics[width=0.8\textwidth]{hydration/first-render.png}
    \caption{Scambio di dati tra server e client durante il primo render}
    \label{fig:ytcsr}
\end{figure}

La prima fase è la renderizzazione del contenuto lato server. In questa fase, il server genera il codice HTML di base per la pagina web, ma a differenza del CSR, la pagina generata non è vuota. Invece, il server genera il codice HTML con il contenuto iniziale della pagina, che nel caso di YouTube è priva di dati veri e propri, ma contiene solo il layout di base della pagina, come mostrato in figura~\ref{fig:ytskeleton}, ma nel caso di pagine come dashboard amministrative, la pagina potrebbe contenere i dati reali iniziali.

\vspace{5mm}
\begin{figure}[h]
    \centering
    \includegraphics[width=0.8\textwidth]{hydration/yt-skeleton.png}
    \caption{Schermata home di YouTube dopo il primo render}
    \label{fig:ytskeleton}
\end{figure}
\vspace{3mm}

A questo punto del processo, la pagina web ha raggiunto quella che si chiama \emph{First Contentful Paint} (FCP), ovvero il momento in cui il browser inizia a visualizzare il contenuto della pagina. Questo è un concetto importante per l'esperienza utente, perchè è il momento in cui l'utente percepisce che la pagina sta caricando, e nonostante la pagina non sia ancora pronta ed interattiva, è comunque meglio della pagina bianca che si ottiene con il CSR.

\subsection{Seconda fase: hydration lato client}

\vspace{5mm}
\begin{figure}
    \centering
    \includegraphics[width=0.8\textwidth]{hydration/second-render.png}
    \caption{Scambio di dati tra server e client durante l'hydration}
    \label{fig:ythydration}
\end{figure}

La seconda fase è la hydration del contenuto lato client. In questa fase, il client esegue il codice JavaScript, inizia
a scaricare i dati necessari per popolare la pagina e ad aggiungere gli event handler, e una volta che i dati sono stati
scaricati, il client aggiorna l'albero DOM. Questo processo è chiamato hydration, perchè il client "idrata" la pagina statica con i dati necessari per renderizzarla completamente e renderla interattiva.
A questo punto, l'utente può interagire con essa come se fosse una pagina web tradizionale. Questo momento è chiamato \emph{time-to-interactive} (TTI), ed è un altro concetto importante per l'esperienza utente, perchè è il momento in cui l'utente può interagire con la pagina.